Base drag is the least-constrained term in slender-body supersonic aerodynamics, yet its
weight in an \emph{integrated} apogee prediction is rarely isolated. We ask how heavily the
base-drag term weighs on the apogee of fin-stabilized supersonic sounding rockets relative
to the other supersonic submodels, and whether the closure that carries that weight can be
anchored to external base-pressure data. The base-drag closure as implemented---turbulent
Chapman--Korst shear-layer theory, a Chapman laminar correlation, an ESDU~77021
boundary-layer-thickness correction, an empirical power-law base-drag correlation, a
Viswanath boattail correction, and a finned-base augmentation term---is benchmarked against
external, out-of-sample base-pressure measurements, and a controlled mechanism ablation then
disables each drag contribution in isolation across an archived 24-flight subset (23
single-stage flights plus the AeroPac~104K two-stage closure) spanning Mach 0.54 to 3.46.
The turbulent closure reproduces the reference NACA~TN~3393 nonlifting-body base pressures
to a mean absolute percentage error of 15.9\% (and the Hart free-flight subset to 4.0\%),
and the Chapman laminar branch to 4.4\%. The ablation shows that base-drag augmentation
moves the mean apogee error by 8.10 percentage points---the dominant mechanism, roughly nine
times the next supersonic submodel (compressible skin friction, 0.87 percentage points) and
about 54 times the shock-geometry pre-pass (0.15 percentage points). A decontaminated
prospective holdout, on which the model is more accurate than on the development flights
(mean absolute error 3.95\% versus 5.47\%), shows that the corpus-frozen base-drag constants
generalize rather than overfit. We frame these results as a mechanism-attribution study,
backed by external base-pressure benchmarks, that locates base drag at the top of the apogee
error budget---not a per-closure intercomparison, and not an integrated-validation or parity
claim---complementing the integrated study reported separately.
